La corrección manual de exámenes escritos sigue siendo, en una parte
significativa de la docencia universitaria, un proceso costoso, lento y
con una trazabilidad limitada. En asignaturas con cientos de estudiantes,
consume decenas de horas del equipo docente por convocatoria; introduce
una variabilidad inevitable entre correctores y, en exámenes con un único
evaluador, sesgos sistemáticos difíciles de detectar; complica la
conservación ordenada de los exámenes a efectos de revisión y auditoría; y
retrasa la publicación de calificaciones, postergando a su vez el
\textit{feedback} que el estudiante necesita para aprender. Cualquier
mejora sustancial en este proceso libera tiempo del profesorado y, al mismo
tiempo, mejora la experiencia del alumnado.

Existen plataformas comerciales que abordan este problema, las cuales se
analizan en detalle en el \cref{CAP:STATE-OF-THE-ART}. Más allá de la
implementación funcional de un sistema de corrección digital, el interés
técnico de este trabajo reside en el conjunto de retos que su construcción
plantea y que se detallan en \ref{SEC:INTRO-OBJECTIVES}. El Sistema de
Corrección Digital de Exámenes Escritos (en adelante, SCDEE) se ha
concebido y construido alrededor de estos retos, que constituyen el núcleo
técnico del presente \acs{tfg}.

La construcción del SCDEE se ha repartido entre dos Trabajos de Fin de
Grado coordinados. El presente trabajo aborda el \dfn{backend} del sistema,
mientras que un \acs{tfg} paralelo, presentado por Alejandro Soler
Sichling, aborda el \dfn{frontend}.
